\subsection*{输入与输出不必住在同一空间}
删掉高度的线性投影 $P(x,y,z)=(x,y)$ 接受三个输入分量，却只给两个输出分量，所以矩阵有两行三列。把它想象成从三维场景保留一个平面视图，比先背矩阵尺寸规则更自然。

$m\times n$ 矩阵的每一列记录一个输入基方向在 $m$ 维输出空间中的去向，一共需要记录 $n$ 根输入基。因此行数是输出维度，列数是输入维度；检查乘法能否进行，也就是检查前一步输出是否能交给后一步。

\subsection*{“撤销”需要说明是哪一边}
每个平面点 $(a,b)$ 都有原像 $(a,b,0)$，所以 $P$ 满射。但 $(a,b,5)$ 也投到同一点，所以它不单射。反向嵌入 $J(a,b)=(a,b,0)$ 只是约定一个高度，并没有找回输入原来的高度。

先嵌入再投影，平面点原样返回，得到 $PJ=I_2$。先投影再嵌入，三维点变成 $(x,y,0)$，所以 $JP\ne I_3$。一对矩阵能在一边互相撤销，却不能在另一边撤销，这就是左右逆要分别讨论的原因。

画几条平行于高度轴的线，把每条线看成一组无法区分的输入。这些线是核的平移。以后遇到数据降维、神经网络宽窄不同的层或欠定方程，先数维度，再问保留了什么、丢失了什么。维度分析不能独自解出问题，却经常能在计算前排除不可能的答案。
